\subsection{FX Market Microstructure}

\begin{principle} \hlt{FX Is Dealer-Centric and Fragmented}\\
Spot FX has no single consolidated exchange. Liquidity is split across dealer streams, ECNs, prime-brokered access, bank internalisers, and bilateral relationships.
\end{principle}

\begin{taxonomy} \hlt{FX Venue Segmentation}
\begin{enumerate}[label=\roman*.]
\setlength{\itemsep}{0pt}
\item \hlt{dealer streams}: bilateral executable prices from banks or non-bank market makers;
\item \hlt{ECNs}: central electronic books with disclosed or anonymous participation;
\item \hlt{RFQ}: request-for-quote negotiation for size or less liquid pairs;
\item \hlt{internalisation}: dealer offsets client flow before external hedging;
\item \hlt{fixing windows}: benchmark-driven flow concentrated around reference-rate times.
\end{enumerate}
\end{taxonomy}

\begin{definition} \hlt{Last Look}\\
Last look gives a liquidity provider a short window to accept or reject an incoming trade request.\\
It reduces stale-quote risk for the provider, but creates execution uncertainty and potential information-leakage concerns for the client.
\end{definition}

\begin{caveat} \hlt{FX Cost Is Relationship-Dependent}\\
The displayed spread is only part of the cost. Credit, prime-broker limits, rejection rates, hold times, skewed streams, and information leakage can dominate fill quality.
\end{caveat}
